% Section "Formal verification" of PAPER_TR26_151_PROBLEM1.tex (same header and claim boundary as the main
% file). Source of every statement in this section: lean_theorem_l/BUILD_RECORD.md, axioms_output.txt,
% axioms_output_layer_ranks.txt, verify_by_main_session.txt and verify_layer_ranks_by_main_session.txt
% (independent re-runs), referee_g10/ (audit of the statements of TheoremL.lean). Source of the paragraph
% "Clean-machine build" (added on 27 November 2026): public_release_tr26_151/lean/records/ci_run_1.md.
Theorem~\ref{thm:L} and Corollaries~\ref{cor:triples} and~\ref{cor:proj} have been formalised in Lean~4 (version
4.30.0-rc2) with Mathlib (tag \texttt{v4.30.0-rc2}), in two files of the folder \texttt{lean\_theorem\_l}:
\texttt{TheoremL.lean} ($1065$ lines, SHA-256 \texttt{3019bee1\,\dots\,6b3cef92}) and
\texttt{TheoremL\_LayerRanks.lean} ($1354$ lines, SHA-256 \texttt{a1b16696\,\dots\,21adeefd}), whose first
$1065$ lines are the first file, byte for byte. The full hashes are in the build record. Neither file contains
\texttt{sorry}, an added axiom or \texttt{native\_decide}. The main declarations are:
\begin{center}\small
\begin{tabular}{@{}p{0.39\textwidth}p{0.58\textwidth}@{}}
\toprule
declaration & statement\\
\midrule
\texttt{core\_bound} & orthonormal $r$-frames $u(p,\cdot)$ in a real space $E$, orthogonal on adjacent
points: $r\,n^2\le5\dim E$\\
\texttt{core\_layer\_ranks}$^{\ast}$ & the same hypotheses: $3\,r\,n^3\le5\sum_\Lambda\rank(\Lambda)$
(Theorem~\ref{thm:L} over $\R$)\\
\texttt{theorem\_L\_real} & unit vectors in $\R^d$: $n^2\le5d$ \ (\texttt{\_nine}: $n^2\le9d$)\\
\texttt{theorem\_L\_complex\_hermitian} & unit vectors in $\C^d$, Hermitian form: $n^2\le5d$\\
\texttt{theorem\_L\_projective} & subspaces of dimension $r$ of a real space: $r\,n^2\le5\dim E$\\
\texttt{theorem\_L\_projective\_complex} & the same over $\C$\\
\texttt{theorem\_L\_projectors} & self-adjoint operators on $\R^D$ with ranges of dimension $r$ and
$P_pP_q=0$ on adjacent points: $r\,n^2\le5D$\\
\texttt{corollary\_L\_real} & non-zero vectors in $\R^d$ indexed by the $3$-subsets of $[t]$, orthogonal
when two subsets share one element: $\floor{t/3}^2\le5d$ \ (\texttt{\_nine}: $\le9d$)\\
\texttt{corollary\_L\_complex\_hermitian}$^{\ast}$ & the same in $\C^d$ with the Hermitian form\\
\texttt{corollary\_L\_projective\_complex}$^{\ast}$ & complex subspaces of dimension $r$ indexed by the
$3$-subsets, orthogonal when two subsets share one element: $r\floor{t/3}^2\le5\dim_\C E$
(Corollary~\ref{cor:proj} for ranges)\\
\texttt{upper\_bound\_triples}$^{\ast}$ & a representation of $\Gt$ by non-zero vectors of $\R^{t^2}$ exists\\
\bottomrule
\end{tabular}
\end{center}
($^{\ast}$: in the second file.) The statement about the quantity of Problem~1 of~\cite{HI} reads, with the
Unicode symbols of Lean transliterated:
\begin{verbatim}
theorem corollary_L_real {t d : Nat}
    (v : {A : Finset (Fin t) // A.card = 3} -> EuclideanSpace Real (Fin d))
    (h_ne : forall A, v A != 0)
    (h_orth : forall A B, (A.1 inter B.1).card = 1 ->
                inner Real (v A) (v B) = 0) :
    (t / 3) ^ 2 <= 5 * d
\end{verbatim}
Here \texttt{t / 3} is the division of natural numbers. Together with \texttt{upper\_bound\_triples} (a
representation by non-zero vectors indexed by $\texttt{Fin}\ t\times\texttt{Fin}\ t$, that is, in dimension
$t^2$) it gives $\floor{t/3}^2/5\le d_\R(t)\le t^2$; the minimum $d_\R(t)$ itself is not defined in the Lean
files.

\paragraph{Checks.} Each file was compiled by the Lean kernel twice, by the process that wrote it and in an
independent re-run, with exit code $0$ and no warning. For every declaration of the table,
\texttt{\#print axioms} reports exactly \texttt{propext}, \texttt{Classical.choice} and \texttt{Quot.sound}.
As a negative control, a file that states \texttt{theorem\_L\_real} and proves it by \texttt{sorry} is
reported as depending on \texttt{sorryAx}. A further theorem shows that the hypotheses of
\texttt{theorem\_L\_real} can be satisfied for every $n$, so that statement is not vacuous.
% Result of the isolated audit of the Lean statements (referee_g10/REFEREE_REPLY.md, 26 November 2026).
The formal statements were then compared with the definitions of~\cite{HI} (its pages 3--4 and~14) by an
isolated reader who did not run Lean. The reader found \texttt{corollary\_L\_real} to state exactly the lower
bound asked for in Problem~1 (vertices: the $3$-subsets; non-zero vectors of $\R^d$ with the standard inner
product; orthogonality required exactly for two subsets that share one element; no further hypothesis), and
found no construct in the file that changes the meaning of a statement. The reader's caveats concern
coverage only and are listed in the next paragraph. A second isolated reader audited the declarations of the
second file in the same way, with the same conclusion.


\paragraph{What is not formalised.} The formal proof follows Section~\ref{sec:proof} in the frame form
($u(p,k)$ in place of $U_p$) and over $\R$; the complex statements are derived inside Lean from the real ones
(a complex subspace of dimension $r$ is a real subspace of dimension $2r$), not by repeating the proof over
$\C$. The following are \emph{not} formalised: the definitions of $\xi_f$ and $\chi_q$ and the inequality
$\xi_f\le\chi_q$ quoted from~\cite{MR}; the passage from complex projectors with $P_AP_B=0$ to orthogonal
ranges (formalised over $\R$ only); Remark~\ref{rem:k}; the statements of Section~\ref{sec:mink}; and all of
Section~\ref{sec:linear} (the certificates of Theorem~\ref{thm:cert} are checked by exact-arithmetic scripts,
not by Lean). The declarations marked $^{\ast}$ were added after the first audit; a second isolated
reader then compared them with the mathematics (\texttt{referee\_g11}) and found them aligned, with caveats
about coverage that are all listed in this paragraph.

\paragraph{Clean-machine build (added on 27 November 2026).} The checks above were run against an existing
Mathlib build, on the machine on which the files were written. The build kit for a clean machine
(\texttt{lakefile.toml}, \texttt{lean-toolchain} and the workflow \texttt{lean\_verify}) was then executed on a
fresh GitHub-hosted runner, by GitHub Actions (run~1 of the workflow, 27~November 2026), and every step
succeeded. The build of the two files printed $27$ lines of \texttt{\#print axioms} output ($10$ from the
first file, $17$ from the second), each naming exactly \texttt{propext}, \texttt{Classical.choice} and
\texttt{Quot.sound}; the planted \texttt{sorry} of the negative control was reported as \texttt{sorryAx}; and
the exact-arithmetic verifier accepted all seven certificates of Theorem~\ref{thm:cert}. The record of the run
is the file \texttt{lean/records/ci\_run\_1.md} of the repository
\url{https://github.com/aqida-ai/tr26-151-problem1}. That run used the same version of Lean and the same
Mathlib as the first checks, and the Mathlib binaries were downloaded from the Mathlib cache, not rebuilt. No
proof checker other than the Lean kernel was used.
